% Note accompanying the mobius-ramsey repository (Lean package in the repository root).
% Engine: LuaLaTeX or XeLaTeX (the Lean statements are input verbatim, in Unicode).
% The quotations and Lean excerpts are checked by note/check_verbatim.py (see REPRODUCE.md).

\documentclass[11pt]{amsart}
\usepackage{iftex}
\ifPDFTeX
  \errmessage{Compile with lualatex or xelatex: the Lean statements are Unicode}
\fi
\usepackage{fontspec}
% JuliaMono covers every glyph in the Lean excerpts (R, N, ^T, _v, arrows, top, circ);
% check glyph coverage at first build if another mono font is substituted.
\IfFontExistsTF{JuliaMono}{\setmonofont{JuliaMono}[Scale=MatchLowercase]}
  {\setmonofont{DejaVu Sans Mono}[Scale=MatchLowercase]}
\usepackage[margin=1in]{geometry}
\usepackage{amsmath,amssymb,amsthm,mathrsfs}
\usepackage{booktabs}
\usepackage{verbatim}
\usepackage{xcolor}
\usepackage[hidelinks]{hyperref}

\newcommand{\TODO}[1]{\textcolor{red}{[TODO: #1]}}
\newcommand{\PH}[1]{\textcolor{blue}{[#1]}}
\newcommand{\st}[1]{\textsf{[#1]}}
\newcommand{\Kriz}{K\v{r}\'{\i}\v{z}}
\newcommand{\Q}{\mathbb{Q}}
\newcommand{\R}{\mathbb{R}}
\newcommand{\C}{\mathbb{C}}
\DeclareMathOperator{\Sym}{Sym}
\DeclareMathOperator{\PO}{PO}
\DeclareMathOperator{\Ort}{O}
\DeclareMathOperator{\PGL}{PGL}
\DeclareMathOperator{\Mob}{M\ddot{o}b}
\DeclareMathOperator{\aff}{aff}
\DeclareMathOperator{\Mat}{Mat}
\newcommand{\circmap}{\mathrm{circ}}

% Lean names (the switch changes the theorem name for Theorem C).
\newcommand{\LeanThmCFullName}{\texttt{OAI.EuclideanRamsey.Mobius.ramsey\_iff\_ramsey\_mobius\_full}}
\newcommand{\LeanCorName}{\texttt{ramsey\_mobius\_of\_algebraic}}
\newcommand{\LeanCorCircName}{\texttt{ramsey\_circ\_mul\_rat}, \texttt{ramsey\_circ\_add\_rat}}

\theoremstyle{plain}
\newtheorem*{theoremC}{Theorem C}
\newtheorem{corollary}{Corollary}
\newtheorem*{oaithm}{Theorem 1.1}
\newtheorem*{oaiprop}{Proposition 2.1}
\theoremstyle{definition}
\newtheorem{remark}{Remark}

\title{Ramsey-ness of finite spherical sets is M\"obius invariant}
\author{Will Strinz}
\date{October 2026. Archived on Zenodo; the DOI is listed in the repository README.}

\begin{document}

\begin{abstract}
OpenAI's classification of finite Euclidean Ramsey configurations \cite{OAI172} decides whether
a finite spherical set is Ramsey by a tensor condition over its coordinate field. We note that
this condition is transported by the Lorentz group: whether a finite subset of a sphere is
Ramsey is invariant under M\"obius transformations of the sphere. On a circle, it therefore
depends only on cross-ratios. Consequences: every M\"obius image of a finite spherical set
with algebraic coordinates is Ramsey; so is every M\"obius image of a regular polygon; and so
are the families $\{\circmap(kt_j)\}$ and $\{\circmap(t_j+c)\}$ in the stereographic
parameter, for distinct rationals $t_j$ and every real $k\ne0$ and $c$, transcendental or not.
The result is a corollary of \cite{OAI172}, and the deep work is OpenAI's.
For the unit sphere, with M\"obius transformations given by Lorentz matrices, the invariance
theorem is proved in Lean~4 on top of OpenAI's formalization, and so are the corollaries for
algebraic coordinates and for the two circle families; the proofs are checked against a
Mathlib-only statement of the theorems by Comparator with three independent kernels. The
passage to an arbitrary sphere, the regular-polygon case and the identification of the Lorentz
model with the M\"obius group are not formalized: the first two are proved by hand, and the
identification is cited.
\end{abstract}

\maketitle

\noindent\textbf{Status conventions.} Each statement carries one tag.
\st{LEAN}: proved in Lean~4 (v4.35.0-rc2, Mathlib \texttt{065356127b}) on top of OpenAI's Lean
proof of \cite{OAI172} (ported to the Lean module system by header-only changes), with
\texttt{\#print axioms} reporting only \texttt{propext}, \texttt{Classical.choice},
\texttt{Quot.sound}.
\st{HAND}: a complete written argument, not machine-checked. The Comparator \cite{Comparator}
check of each \st{LEAN} item is reported in \S\ref{sec:lean}. Each statement
also names the type of its coordinates over $\Q$ (\texttt{ANY}, \texttt{ALGEBRAIC},
\texttt{MOBIUS\_ORBIT\_OF(ALGEBRAIC)}). ``Real'' alone is never a scope, because Ramsey-ness can
differ between algebraic and transcendental members of one real family
\cite[Prop.~7.6]{OAI172}.

\section{The classification}\label{sec:oai}

A finite set $X$ in a Euclidean space is \emph{Ramsey} if for every $r\ge2$ there is $D$ such that
every $r$-colouring of $\R^D$ is constant on some congruent copy of $X$
\cite{EGMRSS,OAI172}. Ramsey sets are spherical \cite{EGMRSS}. Graham's conjecture, that every
spherical set is Ramsey, was refuted by P\'alv\"olgyi \cite{Palvolgyi}. OpenAI \cite{OAI172}
then classified the finite Ramsey configurations. We quote the two statements we use verbatim
from the \TeX{} source at commit \texttt{fd4aeeb2} (Apache-2.0). The only edits are cross-reference
labels, which we print as OpenAI's rendered numbers. From \cite[\S1.1]{OAI172}:

{\small
\begin{quote}
A singleton is Ramsey. For any other set, let $d\ge1$ be its affine
dimension, and choose a congruent representative
$A=\{a_1,\ldots,a_s\}\subset\R^d$ whose affine span is $\R^d$.
The points are distinct and $s\ge2$. Define
\[
 p_i=\begin{pmatrix}1\\a_i\end{pmatrix},\qquad
 F=\Q\bigl((a_i)_\alpha:1\le i\le s,\ 1\le\alpha\le d\bigr),
 \qquad B=F\otimes_{\Q}F.
\]
The commutative ring $B$ has a multiplication homomorphism
\[
 m_F:B\longrightarrow F,\qquad m_F(x\otimes y)=xy.
\]
An equality in $B$ is stronger than the equality obtained by applying
$m_F$; this map need not be injective. Index the entries of $p_i$ and
the rows and columns of a $(d+1)$-by-$(d+1)$ matrix by $0,1,\ldots,d$.
Thus index $0$ is the constant coordinate.

\begin{oaithm}[Classification]
Let $A=\{a_1,\ldots,a_s\}\subset\R^d$ be a finite set of at least two
points with affine span $\R^d$, and define $p_i,F,B,m_F$ as above.
Then $A$ is Ramsey if and only if there is a matrix
$P\in\operatorname{Mat}_{d+1}(B)$ such that
\begin{align}
 (p_i\otimes1)^T P(1\otimes p_i)&=0
       &&(1\le i\le s),\tag{1}\\
 m_F(P_{\alpha\beta})&=\delta_{\alpha\beta}
       &&(1\le\alpha,\beta\le d).\tag{2}
\end{align}
Here the entries of $p_i\otimes1$ and $1\otimes p_i$ are formed
componentwise. There are no separate conditions on the multiplied
constant row or constant column of $P$.
\end{oaithm}
\end{quote}
}

\noindent From \cite[\S2]{OAI172}:

{\small
\begin{quote}
Throughout, $s\geq2$, $d\geq1$, and
$A=\{a_1,\ldots,a_s\}$ affinely spans $\R^d$; the field $F$, ring $B$,
columns $p_i$, and multiplication map $m_F$ are as in the theorem.

Put $W=\R\times\R^d$, regarded as a vector space over $\Q$.
An element $w=(v,u)$ specifies the affine function $a\mapsto v+a\cdot u$.
For each label define the $\Q$-linear evaluation map
\[
 e_i:W\longrightarrow\R,\qquad e_i(v,u)=v+a_i\cdot u.
\]
Let $\mathcal T(W)$ be the subspace of $W\otimes_{\Q}W$ fixed by the
flip $w\otimes w'\mapsto w'\otimes w$.  Thus our symmetric tensors are
flip-invariant tensors.  The map
\[
 (v,u)\otimes(v',u')\longmapsto u(u')^{\mathsf T}
\]
is $\Q$-linear on the tensor product and restricts to a map
$G:\mathcal T(W)\to\Sym_d(\R)$, where $\Sym_d(\R)$ denotes the real
symmetric matrices.  In particular, $G(w\otimes w)=uu^{\mathsf T}$.
The tensor condition is the existence of $T\in\mathcal T(W)$ such that
\begin{equation}
 (e_i\otimes e_i)(T)=0\quad(1\leq i\leq s),
 \qquad G(T)=I_d.\tag{3}
\end{equation}
The first identities hold in $\R\otimes_{\Q}\R$.
They are stronger than their images under real multiplication.

\begin{oaiprop}
The matrix condition in Theorem~1.1 is equivalent
to~(3).
\end{oaiprop}
\end{quote}
}

\noindent\emph{Array form.} Write $\mathscr B=\R\otimes_\Q\R$, $m:\mathscr B\to\R$ for
multiplication and $\tau$ for the flip. The proof of Proposition~2.1 identifies
$W\otimes_\Q W$ with $(d+1)\times(d+1)$ arrays $Q$ over $\mathscr B$. Under this identification
$(e_i\otimes e_i)(T)$ becomes $(p_i\otimes1)^TQ(1\otimes p_i)$, $G(T)$ becomes the spatial block of
$m(Q)$, and the flip becomes $Q^*_{\alpha\beta}=\tau(Q_{\beta\alpha})$. Since $m\circ\tau=m$, we
have $m(Q^*)=m(Q)^T$. That proof also shows that $\tfrac12(Q+Q^*)$ keeps every vanishing
evaluation and a multiplied spatial block equal to $I_d$. We use (3) in this form, with
$B\subset\mathscr B$. In OpenAI's Lean proof, the array form is \texttt{RealCriterion}; the
necessity and descent halves we use are \texttt{ramsey\_realCriterion} and
\texttt{realCriterion\_fieldCriterion}, and sufficiency is \texttt{fieldCriterion\_ramsey}.

\section{M\"obius invariance}\label{sec:mobius}

Let $S=\mathrm{diag}(-1,1,\dots,1)$ on $\R^{1+d}$, with coordinates indexed $0,1,\dots,d$ as above,
and let $\Ort(1,d)=\{L\in\Mat_{d+1}(\R):L^TSL=S\}$. For $x$ on the unit sphere
$S^{d-1}\subset\R^d$, the vector $(1,x)$ is $S$-null. If $L\in\Ort(1,d)$, then $L(1,x)$ is again
null and nonzero, so its $0$-th coordinate $\lambda$ is nonzero: a null vector with vanishing
$0$-th coordinate is $0$. Put
\[
 L\begin{pmatrix}1\\x\end{pmatrix}=\lambda\begin{pmatrix}1\\\varphi_L(x)\end{pmatrix},
 \qquad \varphi_L(x)\in S^{d-1}.
\]
For $d\ge2$, the maps $\varphi_L$ are the M\"obius transformations of $S^{d-1}$. Here the M\"obius
group $\Mob(S^{d-1})$ consists of the restrictions to $S^{d-1}$ of the M\"obius transformations of
$\R^d\cup\{\infty\}$ that preserve $S^{d-1}$. For $d\ge3$ these are exactly the conformal
diffeomorphisms of $S^{d-1}$ (Liouville). For $d=2$ they are $\PGL_2(\R)$ acting on the
stereographic parameter. For $d\ge2$, the map $L\mapsto\varphi_L$ induces the classical isomorphism
$\PO(1,d)=\Ort(1,d)/\{\pm I\}\cong\Mob(S^{d-1})$, orientation-reversing maps included
\cite[Thm.~3.2.3 and \S3.2, Cor.~3; \S4.4, Cor.~3; Thm.~4.5.2 and \S4.5, Cor.~1]{Ratcliffe},
\cite[\S3.7, Thms.~3.7.5 and 3.7.6, pp.~49--54]{Beardon}. In Ratcliffe's notation it is the
composite of Poincar\'e extension $\Mob(S^{d-1})\cong M(B^d)$, restriction $M(B^d)\cong I(B^d)$,
the isometry of the conformal ball model $B^d$ with $H^d$, and $I(H^d)\cong\Ort^+(d,1)$; his time
coordinate is the last one. Beardon puts the time coordinate first,
as here, and acts on row vectors; he gives the isomorphism of $\Ort^+(1,d)$ onto the full M\"obius
group of $B^d$ by conjugation with his map $F$ from the hyperboloid to the ball.
We cite this identification. It is not formalized.
The restriction $d\ge2$ matters: for $d=1$ the sphere $S^0=\{\pm1\}$ has two points, every boost
in $\Ort(1,1)$ fixes both, and the maps $\varphi_L$ are just the two permutations of $S^0$, so
$L\mapsto\varphi_L$ is not injective on $\PO(1,1)$. Theorem~C still holds for $d=1$ (every
subset of $S^0$ is Ramsey: a singleton is, and every two-point set is \cite[Cor.~7.4]{OAI172}),
and the Lean statement covers it, but there it is not a statement about a M\"obius group. For
$d=0$ the unit sphere is empty.

\begin{theoremC}
Let $A$ be a finite subset of a sphere in $\R^d$, and let $\varphi$ be a M\"obius transformation
of that sphere. Then $A$ is Ramsey if and only if $\varphi(A)$ is Ramsey. In particular, for
points on a circle, Ramsey-ness depends only on cross-ratios (Corollary~\ref{cor}(a)).

\smallskip\noindent\emph{Status:} \st{LEAN} for $s$ distinct points of the unit sphere and
$\varphi=\varphi_L$, $L\in\Ort(1,d)$, with no spanning hypothesis. Lean: \LeanThmCFullName{}
(\S\ref{sec:lean}), accepted by Comparator. The reduction from an arbitrary sphere to the unit
sphere is a similarity \st{HAND} (Remark~\ref{rem:wlog}), and the identification of
$\Mob(S^{d-1})$ with $\PO(1,d)$ is cited. Scope: \texttt{ANY}.
\end{theoremC}

\begin{proof}[Proof of Theorem C for spanning sets]
Write $A=\{a_i\}$, $A'=\{\varphi_L(a_i)\}$, $p_i=(1,a_i)$ and $p_i'=(1,\varphi_L(a_i))$, so that
$Lp_i=\mu_ip_i'$ with $\mu_i\in\R^\times$. A singleton is Ramsey, so assume $s\ge2$; then $d\ge1$.

\emph{Step 1: $A'$ is again a configuration of the same kind.} If $x,y\in S^{d-1}$ are distinct,
then $(1,x)$ and $(1,y)$ span distinct lines, and so do their images under the invertible $L$.
Hence $\varphi_L(x)\neq\varphi_L(y)$, and the $s$ points of $A'$ are distinct and lie on $S^{d-1}$.
A finite set affinely spans $\R^d$ if and only if its augmented vectors span $\R^{1+d}$, that is,
if and only if no nonzero affine function vanishes on it. The vectors $p_i'=\mu_i^{-1}Lp_i$ span
$\R^{1+d}$ because the $p_i$ do. So $A'$ affinely spans $\R^d$.

\emph{Step 2: a normalized certificate for $A$.} Suppose $A$ is Ramsey. Theorem~1.1 gives $P$ over
$B\subset\mathscr B$ with (1) and (2). Replace $P$ by $\tfrac12(P+P^*)$. This keeps (1), and
$m(P)$ becomes symmetric with spatial block $I_d$. Applying $m$ to (1) gives
$(1,a_i)\,m(P)\,(1,a_i)^T=0$, so the quadratic function $q(x)=(1,x)\,m(P)\,(1,x)^T$ vanishes on $A$.
It has the form $\|x\|^2+\ell\cdot x+c$. So does $(1,x)S(1,x)^T=\|x\|^2-1$, which also vanishes on
$A$. Their difference is an affine function vanishing on the spanning set $A$, hence zero. Thus
$m(P)=S$.

\emph{Step 3: transport.} Let $N=L^{-1}$. Write $N\otimes1$ and $1\otimes N$ for the matrices with
entries $N_{\alpha\beta}\otimes1$ and $1\otimes N_{\alpha\beta}$ in $\mathscr B$, and put
\[
 P'=(N\otimes1)^T\,P\,(1\otimes N).
\]
Since $Np_i'=\mu_i^{-1}p_i$, we get
\[
 (p_i'\otimes1)^TP'(1\otimes p_i')
 =(\mu_i^{-1}\otimes\mu_i^{-1})\,(p_i\otimes1)^TP(1\otimes p_i)=0.
\]
Multiplication is a ring map, so $m(P')=N^TSN=S$, because $L^TSL=S$ implies $N^TSN=S$. The
spatial block of $S$ is $I_d$. If $P^*=P$, then $P'^*=P'$, by applying $\tau$ entrywise.

\emph{Step 4: descent and conclusion.} In the array form, $P'$ is a flip-invariant tensor that
satisfies (3) for $A'$. By Step~1, $A'$ meets the standing hypotheses of \cite[\S2]{OAI172}.
Proposition~2.1 for $A'$ therefore gives a matrix over $B'=F'\otimes_\Q F'$, where $F'$ is the
coordinate field of $A'$, satisfying (1) and (2) for $A'$. Theorem~1.1 then shows that $A'$ is
Ramsey.

The descent step carries the weight. The entries of $L$ in general lie in neither $F$ nor $F'$,
so the transport has to take place over $\mathscr B$.

\emph{Converse.} $L^{-1}=SL^TS\in\Ort(1,d)$ and $\varphi_{L^{-1}}\circ\varphi_L=\mathrm{id}$ on
$S^{d-1}$. Since $A'$ spans by Step~1, the same argument applied to $A'$ and $L^{-1}$ shows that
$A'$ Ramsey implies $A$ Ramsey.
\end{proof}

The Lean proof follows these steps, with the spanning hypothesis removed as described at the end
of Remark~\ref{rem:wlog}. The matrix identities of Step~3 are proved over an arbitrary
commutative ring with two ring maps from $\R$. OpenAI's lemmas named in \S\ref{sec:oai} are
imported, with their statements and proofs unchanged. The project owner read this proof line by
line and ratified it. He describes himself as not an expert in the area, so that read is an owner's sign-off, not expert review.

\begin{remark}[The spanning hypothesis]\label{rem:wlog}
\st{HAND}
Let $A$ be a finite subset of a sphere $\Sigma_0\subset\R^d$ with $|A|\ge2$, and let $\varphi$ be
a M\"obius transformation of $\Sigma_0$. Ramsey-ness is invariant under similarities. Congruences
preserve it by definition. If an $r$-colouring $c$ of $\R^D$ is given, then
$x\mapsto c(\lambda x)$ is another $r$-colouring, and a monochromatic copy of $X$ for the second
colouring gives a monochromatic copy of $\lambda X$ for $c$. So we may take $\Sigma_0=S^{d-1}$.

Let $E=\aff(A)$, of dimension $k$, and let $\Sigma=S^{d-1}\cap E$. This is the smallest subsphere
containing $A$, a $(k-1)$-sphere of positive radius, and $A$ affinely spans $E$. M\"obius maps send
subspheres to subspheres, so $\Sigma'=\varphi(\Sigma)$ is a $(k-1)$-subsphere, with affine hull
$E'$. Choose similarities $\sigma:E\to\R^k$ and $\sigma':E'\to\R^k$ that carry $\Sigma$ and
$\Sigma'$ onto $S^{k-1}$. Then $\psi=\sigma'\circ\varphi\circ\sigma^{-1}$ is a M\"obius
transformation of $S^{k-1}$. In the Lorentz model, $\varphi=\varphi_L$; $L$ maps the
$(k+1)$-dimensional subspace $U=\operatorname{span}\{(1,a):a\in\Sigma\}$, on which $S$ has
signature $(1,k)$, onto $LU$, and in $S$-orthonormal bases of $U$ and $LU$ the induced map lies in
$\Ort(1,k)$. Theorem~C applies to $\sigma(A)\subset S^{k-1}$ and $\psi$, and both similarities
preserve Ramsey-ness. If $k=1$, then $|A|=2$, and every two-point set is Ramsey
\cite[Cor.~7.4]{OAI172}.
The Lean proof of the span-free statement does not follow this route; nothing
is restricted to a subsphere. It removes the spanning hypothesis at the level of OpenAI's real
criterion, in two places. For necessity, OpenAI's \texttt{full\_affine\_representative} gives a
congruent spanning copy $B\subset\R^k$ of $A$, so that $A=t+Q(B)$ for a linear isometry
$Q:\R^k\to\R^d$; the criterion for $B$ is transferred to $A$ in $A$'s own coordinates by adding
rank-one arrays built from the projection onto $Q(\R^k)^\perp$
(\texttt{realCriterion\_of\_isometry}). For the normalization of Step~2, without spanning the
symmetrized multiplied matrix is $S+e_0w^T+we_0^T$ for an affine function $w$ vanishing on $A$,
and subtracting the rank-one arrays $e_0\otimes w+w\otimes e_0$ (zero evaluations) leaves $S$
(\texttt{exists\_lorentz\_normal}). Step~3 is unchanged, and in Step~4 OpenAI's Lean descent
(\texttt{realCriterion\_fieldCriterion}) and sufficiency (\texttt{fieldCriterion\_ramsey}) carry
no spanning hypothesis, so they apply to $A'$ directly.
\end{remark}

\begin{corollary}\label{cor}
\begin{enumerate}
\item[(a)] \st{HAND}, scope \texttt{ANY}. Let $(a_1,\dots,a_s)$ and $(b_1,\dots,b_s)$, $s\ge4$, be
ordered tuples of distinct points, each tuple lying on a circle in some Euclidean space. Identify
each circle with the unit circle of $\C$ by a similarity. If the cross-ratios satisfy
$[a_1,a_2;a_3,a_i]=[b_1,b_2;b_3,b_i]$ for $4\le i\le s$, then $\{a_i\}$ is Ramsey if and only if
$\{b_i\}$ is. (The cross-ratios of concyclic points are real, so they do not depend on the choice
of similarity.)
\item[(b)] \st{HAND}, scope \texttt{ANY}. Every triangle is Ramsey: three distinct points of a
circle are M\"obius-equivalent to an equilateral triangle, which is Ramsey (it is transitive
\cite[Cor.~7.2]{OAI172}). This agrees with Frankl and R\"odl \cite{FR86}. It is \emph{not} an
independent proof, because it rests on \cite{OAI172}, whose Cor.~7.4 already gives it directly.
\item[(c)] \st{LEAN} for the unit sphere (\LeanCorName, accepted by Comparator;
the passage to an arbitrary sphere and the polygon case are \st{HAND}),
scope \texttt{MOBIUS\_ORBIT\_OF(ALGEBRAIC)}. Every M\"obius image of a finite subset of a sphere
whose coordinates are real algebraic numbers is Ramsey. In particular, every M\"obius image of
the vertex set of a regular polygon is Ramsey.
\item[(d)] \st{LEAN} (\LeanCorCircName, accepted by Comparator),
scope \texttt{MOBIUS\_ORBIT\_OF(ALGEBRAIC)}. Let
$\circmap(t)=\bigl(\tfrac{1-t^2}{1+t^2},\tfrac{2t}{1+t^2}\bigr)$, the unit circle in the
stereographic parameter from $(-1,0)$, as in \cite[Prop.~7.6]{OAI172}. Let $t_1,\dots,t_s$ be
distinct rationals. Then for every real $k\neq0$ the set $\{\circmap(kt_j)\}_j$ is Ramsey, and
for every real $c$ the set $\{\circmap(t_j+c)\}_j$ is Ramsey.
\end{enumerate}
\end{corollary}

\begin{proof}
(a) After similarities, both tuples lie on the unit circle. The M\"obius transformation $\mu$ of
$\C\cup\{\infty\}$ with $\mu(a_j)=b_j$ for $j=1,2,3$ maps the circle through $a_1,a_2,a_3$ onto
the circle through $b_1,b_2,b_3$, so it preserves the unit circle. Since $\mu$ preserves
cross-ratios, $\mu(a_i)=b_i$ for all $i$. Three distinct concyclic points span the plane, so
Theorem~C applies in its spanning form, with $\mu|_{S^1}=\varphi_L$.

(b) The case $s=3$ of the argument in (a).

(c) A finite spherical set with real algebraic coordinates is Ramsey (``Theorem~A''; prior work:
\cite{JDJones,Rogerhu}, and end to end in Lean, for configurations with full affine span,
\cite{Atimics}; see \S\ref{sec:credit}). This repository proves it again in Lean on OpenAI's
development (\texttt{ramsey\_of\_algebraic}), with no spanning hypothesis. Apply Theorem~C. The
Lean form of Theorem~A holds on any sphere; for a sphere other than the unit sphere, Theorem~C
itself is applied after the similarity reduction of Remark~\ref{rem:wlog}, which is the
\st{HAND} step. A regular polygon is similar (not in general
congruent) to $\{e^{2\pi i j/n}\}\subset S^1$, which has algebraic coordinates and spans the plane
when $n\ge3$; similarities preserve Ramsey-ness (Remark~\ref{rem:wlog}).

(d) In the stereographic parameter, $u\mapsto ku$ for $k\neq0$ and $u\mapsto u+c$ are in
$\PGL_2(\R)=\Mob(S^1)$; in the Lean proof they are given by explicit matrices in $\Ort(1,2)$
(\texttt{exists\_mobius\_circ\_mul}, \texttt{exists\_mobius\_circ\_add}). They map the rational
points $\circmap(t_j)$ to the given sets. Apply (c). The Lean theorems \texttt{ramsey\_circ\_mul}
and \texttt{ramsey\_circ\_add} allow algebraic $t_j$.
\end{proof}

For transcendental $k$ or $c$ and $s\ge2$ the sets in (d) have a transcendental coordinate field
($\Q(k)$ or $\Q(c)$: the parameter $u$ of $\circmap(u)=(x,y)$ is $y/(1+x)$, and some $t_j\ne0$),
yet they are Ramsey for every number of points. By contrast, nine points
$\circmap(t_1),\dots,\circmap(t_9)$ with algebraically independent $t_i$ are not Ramsey
\cite[Prop.~7.6]{OAI172}. So for transcendental coordinates, transcendence alone does not decide
the question. Here the M\"obius orbit of a
rational set does.

\section{The Lean statement and its verification}\label{sec:lean}

OpenAI's definitions, used verbatim (\texttt{Model.lean} at \texttt{fd4aeeb2}, Apache-2.0):
{\small\verbatiminput{excerpts/OAI-Model-L9-17.txt}}
\noindent Statement definitions (\texttt{MobiusRamsey/Defs.lean}, lines 32--43 and 48--49;
nine code lines):
{\small\verbatiminput{excerpts/Defs-L32-43.txt}}
{\small\verbatiminput{excerpts/Defs-L48-49.txt}}

\noindent Theorem C (\texttt{MobiusRamsey/Theorem.lean}, lines 68--72; proof omitted):
{\small\verbatiminput{excerpts/Theorem-L68-72.txt}}
\noindent Corollary~\ref{cor}(c), unit-sphere form, and (d)
(\texttt{MobiusRamsey/Corollaries.lean}, lines 30--35, 67--68 and 72--73; proofs omitted):
{\small\verbatiminput{excerpts/Corollaries-L30-35.txt}}
{\small\verbatiminput{excerpts/Corollaries-L67-68.txt}}
{\small\verbatiminput{excerpts/Corollaries-L72-73.txt}}

\noindent\emph{Fidelity.} ``M\"obius transformation of $S^{d-1}$'' is modelled as
\texttt{mobiusMap L} with $L^TSL=S$. The identification with $\Mob(S^{d-1})$ is cited
(\S\ref{sec:mobius}). Off the sphere, \texttt{mobiusMap} is not meaningful (Lean's $x/0=0$), but
the theorems evaluate it only on the sphere, where the divisor is proved nonzero. Points are an
injective family on the unit sphere centred at $0$; a general sphere reduces to it by a
similarity. In Corollary~\ref{cor}(d) the Lean names match the text: \texttt{t i} is the
rational $t_i$, and \texttt{k}, \texttt{c} are the real factor and shift.

OpenAI's \texttt{Ramsey} is used unchanged.
\texttt{\#print axioms} gives \texttt{[propext, Classical.choice, Quot.sound]} for every theorem
quoted here. The statements need 9 lines of definitions beyond Mathlib and OpenAI. The proof of
Theorem~C takes 672 lines (\texttt{Defs}, \texttt{Transport}, \texttt{FullTransport},
\texttt{Theorem}); the corollaries take 367 more (\texttt{Circle}, \texttt{Corollaries}, and
\texttt{AlgebraicRamsey} for Theorem~A).

\smallskip\noindent\emph{Verification.} The repository is a single Lake package. OpenAI's 28
files of the import closure of \texttt{Main} are included with their statements and proofs
unchanged; to meet the Lean module system they were ported by header-only changes, listed in
\texttt{docs/PORT-LOG.md}. \texttt{Challenge.lean} restates the definitions above and the four
theorem statements using Mathlib alone. The toolchain's \texttt{lake comparator} \cite{Comparator}
checks that the proved theorems have exactly these statements and definitions and only the three
standard axioms, and replays the proofs through three kernels (Lean's, NanoDa and con-ron). It
accepts. Five altered challenges (negative controls) are refused for the expected reasons: the
challenge used as its own solution (\texttt{sorryAx}), \texttt{Ramsey} with $3\le r$,
\texttt{lorentzForm} replaced by the identity, a changed \texttt{circ}, and Theorem~C without
injectivity. The last is a stronger statement, probably true; it is refused because it is not the
statement that was proved. The proofs trust
OpenAI's Lean proof of \cite{OAI172}. Repository: \url{https://github.com/wstrinz/mobius-ramsey}; archived on Zenodo (DOI in the repository README); replay: \texttt{REPRODUCE.md} in the repository.

\section{Credit}\label{sec:credit}

\begin{itemize}
\item \textbf{OpenAI \cite{OAI172}} is the foundation: Theorem~1.1, Proposition~2.1, and the Lean
proof on which ours builds. Theorem~C is a short consequence of that work.
\item \textbf{P\'alv\"olgyi \cite{Palvolgyi}} gave the seven-point non-Ramsey heptagon. By the
arXiv comments and the disclosure in the paper itself, the whole manuscript was written by
ChatGPT, and its appendix, which he labels as unchecked, consists of further results by ChatGPT
that he has not verified. The appendix proves in Lemma~A.5 that the
\emph{existence of his derivation certificates} is invariant under real M\"obius maps. It states
explicitly that the lemma ``does not assert that the Ramsey property itself is M\"obius
invariant.'' Theorem~C is that stronger assertion, obtained from \cite{OAI172}. The appendix's
Theorem~A.10 is the certificate-level analogue, on the circle, of Corollary~\ref{cor}(c): a
finite subset of the unit circle that a M\"obius transformation sends to points with algebraic
coordinates, apart from at most three, admits no certificate of his (non-Ramsey) derivation
criterion. The appendix notes that this is a limitation of the certificate method and does not
decide Ramsey-ness.
\item \textbf{JD-Jones-ASES/six-concyclic-lean} \cite{JDJones} (commit \texttt{0703d7d},
2026-10-07) and \textbf{Rogerhu12/openai-math-extensions} \cite{Rogerhu} (commit
\texttt{f8a2bb2}, 2026-10-07) observed that algebraic coordinates satisfy OpenAI's criterion:
JD-Jones in Lean, at criterion level, with the Ramsey conclusion conditional on a hypothesis
\texttt{hsuff}; Rogerhu12 in \TeX. \textbf{atimics/math2} \cite{Atimics} (commit \texttt{8ab51fa},
2026-10-08) proved Theorem~A end to end in Lean on OpenAI's proof, for algebraic configurations
with full affine span (its Lean statement carries a spanning hypothesis); its written derivation
extends this to sets similar to algebraic ones. Corollary~\ref{cor}(c) rests on Theorem~A; this
repository's \texttt{ramsey\_of\_algebraic} is a further Lean proof of it, not a first one.
\item Literature and repository searches found no prior statement of Theorem~C; the search log
is \texttt{docs/PRIOR-WORK-SEARCH.md} in the repository. Leader--Russell--Walters \cite{LRW12}
contains no M\"obius or cross-ratio remark, and neither does \Kriz{} \cite{Kriz91}; we could not
obtain \cite{Kriz92}. Gibbons and Luo \cite{GibbonsLuo} proposed an ``inversive Ramsey theory'',
with the M\"obius group in place of the Euclidean group, but do not relate the Euclidean Ramsey
property to M\"obius images. The searches were last repeated on 2026-10-13. The corollaries in
Corollary~\ref{cor} are derivable by an expert from Theorem~C and known results, and are not
claimed as independently new.
\end{itemize}

\section{AI use and review disclosure}\label{sec:disclosure}

The author is the only human contributor; \texttt{DISCLOSURE.md} in the repository gives the
details. Theorem~C and its proof came out of a conversation with Claude Opus~5.5 (Anthropic,
in claude.ai). Claude Opus~5.5 then coordinated the work, wrote the Lean
proofs, checked the hand proof step by step, and condensed this note from a longer draft; Claude
Sonnet~5.5 carried out the literature searches and other named checks. All reviews of this work
other than the author's read are AI reviews. The author read the hand proof of Theorem~C line by
line; he is not an expert in the area, so that read is an owner's sign-off, not expert review.
\textbf{No human expert has reviewed this note.} Machine checks: Lean and Comparator, for
Theorem~C and for the Lean forms of Corollary~\ref{cor}(c), (d) (\S\ref{sec:lean}).

\begin{thebibliography}{99}
\bibitem{OAI172} OpenAI, \emph{A classification of finite Euclidean Ramsey configurations},
OpenAI Math Release preprint, September 23, 2026; github.com/openai/math, commit
\texttt{fd4aeeb2ee4fc729c18d98444fed42fd0529eeeb}, \texttt{preprints/A-classification-of-finite-Euclidean-Ramsey-configurations-September-23-2026/}.
\bibitem{Palvolgyi} D.~P\'alv\"olgyi, \emph{A cyclic non-Ramsey heptagon}, arXiv:2609.23327 (v1,
2026-09-20; no later version as of 2026-10-13).
\bibitem{JDJones} JD-Jones-ASES, \emph{six-concyclic-lean}, GitHub repository,
\url{https://github.com/JD-Jones-ASES/six-concyclic-lean}, root commit \texttt{0703d7d}, 2026-10-07.
\bibitem{Atimics} atimics, \emph{math2}, GitHub repository, \url{https://github.com/atimics/math2},
\texttt{EndToEnd/AlgebraicRamsey.lean}, commit \texttt{8ab51fa}, 2026-10-08.
\bibitem{Rogerhu} Rogerhu12, \emph{openai-math-extensions},
\texttt{papers/172-euclidean-ramsey/algebraic\_spherical\_configurations.tex}, commit
\texttt{f8a2bb2}, 2026-10-07, \url{https://github.com/Rogerhu12/openai-math-extensions}.
\bibitem{EGMRSS} P.~Erd\H{o}s, R.~L.~Graham, P.~Montgomery, B.~L.~Rothschild, J.~Spencer,
E.~G.~Straus, \emph{Euclidean Ramsey theorems I}, J. Combin. Theory Ser. A \textbf{14} (1973),
no.~3, 341--363. doi:10.1016/0097-3165(73)90011-3.
\bibitem{FR86} P.~Frankl, V.~R\"odl, \emph{All triangles are Ramsey}, Trans. Amer. Math. Soc.
\textbf{297} (1986), no.~2, 777--779. doi:10.1090/S0002-9947-1986-0854099-6.
\bibitem{Kriz91} I.~\Kriz, \emph{Permutation groups in Euclidean Ramsey theory}, Proc. Amer.
Math. Soc. \textbf{112} (1991), no.~3, 899--907. doi:10.1090/S0002-9939-1991-1065087-9.
\bibitem{Kriz92} I.~\Kriz, \emph{All trapezoids are Ramsey}, Discrete Math. \textbf{108} (1992),
59--62. doi:10.1016/0012-365X(92)90660-8.
\bibitem{GibbonsLuo} J.~C.~Gibbons, Y.~Luo, \emph{Colorings of the $n$-sphere and inversive geometry},
arXiv:1307.5351 (2013).
\bibitem{LRW12} I.~Leader, P.~A.~Russell, M.~Walters, \emph{Transitive sets in Euclidean Ramsey
theory}, J. Combin. Theory Ser. A \textbf{119} (2012), no.~2, 382--396.
doi:10.1016/j.jcta.2011.09.005.
\bibitem{Ratcliffe} J.~G.~Ratcliffe, \emph{Foundations of Hyperbolic Manifolds}, 3rd ed.,
Graduate Texts in Mathematics \textbf{149}, Springer, Cham, 2019. doi:10.1007/978-3-030-31597-9.
\bibitem{Beardon} A.~F.~Beardon, \emph{The Geometry of Discrete Groups}, Graduate Texts in
Mathematics \textbf{91}, Springer, New York, 1983. doi:10.1007/978-1-4612-1146-4.
\bibitem{Comparator} leanprover/comparator, as bundled in the Lean toolchain
\texttt{leanprover/lean4:v4.35.0-rc2} (\texttt{lake comparator}, with the NanoDa and con-ron
kernels).
\end{thebibliography}

\end{document}
